% Set from report/draft.md, part 1. The prose is the maintainer's; this file is
% its transit into LaTeX and nothing in it was rewritten.
\section{Introduction}

\subsection{Motivation: active vision as sampling under a budget}

Open a camera's shutter and every pixel arrives at the same resolution. Open an
eye and it does not. The human fovea subtends one to two degrees of high acuity;
beyond it resolution falls off sharply, and three or four times a second a
saccade repositions that window somewhere new. The retina is not a uniform
sensor that attention is later applied to. It is a variable-resolution sensor,
and foveation is a bandwidth decision before it is an attention mechanism.

That distinction is easy to state and, as this report will show, easy to lose in
implementation. It is the hinge on which most of what follows turns.

Binocular vision adds a second sense in which biological vision is \emph{active}.
The eyes are not two cameras with fixed extrinsics. They rotate about centres
fixed in the head; vergence sets the distance at which the two lines of sight
intersect; and the torsion of each eye is not free but determined by gaze
through Listing's law --- the eye's orientation lies in a plane that tilts with
vergence. Stereo geometry in a biological system is therefore \emph{oculomotor}:
it follows from where the system is looking, and it changes every time it looks
somewhere else.

This report describes a project that took that framing literally. The claim
under test is the one the active-perception tradition has advanced since the
1980s --- Bajcsy's \emph{Active Perception}~\cite{bajcsy1988active}, Aloimonos,
Weiss and Bandyopadhyay's \emph{Active Vision}~\cite{aloimonos1988active},
Ballard's \emph{Animate Vision}~\cite{ballard1991animate} --- that perception is
better posed as a controlled sampling problem than as a reconstruction problem.
Applied to stereo depth, the specific question becomes: does choosing where to
look improve depth estimation, relative to not choosing?

The stimulus deserves its own mention. Julesz's random-dot
stereograms~\cite{julesz1960binocular} isolate stereopsis by removing every
monocular depth cue --- depth exists only in disparity, so a system that appears
to recover it cannot be exploiting shading or texture gradient. That property
made them the natural synthetic fixture here. It also made a defect in that
fixture unusually damaging, in a way \S2 takes up.

\subsection{Background: stereo depth as inference}

Stereo matching admits a Bayesian reading. For each pixel, comparing a patch
against candidate positions in the other image produces a cost curve over
disparity --- a likelihood. A prior over depth or a smoothness assumption
supplies the rest, and the quantity of interest is the posterior: an estimate
with an uncertainty attached, not a point estimate.

Uncertainty is usually read from the cost curve. Curvature at the minimum gives
a local precision; the ratio between the best and second-best minima measures
ambiguity; agreement between left-to-right and right-to-left matching flags
pixels with no consistent correspondent. Estimates from multiple views are then
combined by precision weighting --- each measurement contributes in proportion
to its confidence.

Depth is not disparity. Recovering metric depth requires knowing the geometry,
and for a converged binocular system it requires knowing the fixation distance.
This is where the oculomotor framing enters the estimator rather than just the
camera model, and it is the origin of one of the report's more portable
findings.

One property of this pipeline must be stated, because it motivated the work and
because the record complicates it. A half-occluded pixel --- visible to one eye,
hidden from the other --- has no true correspondent. The matcher does not report
this; it finds the best available wrong match, and a sharp cost minimum at a
wrong disparity is genuinely sharp. The result is a confident, incorrect
estimate at exactly the geometry an active system might hope to resolve by
looking elsewhere.

The predecessor repositories measured this on photographs, across three
uncertainty readouts, and found variance in half-occluded regions \emph{lower}
than in matched ones --- confidence inverted where it matters most. On
random-dot stereograms the same matcher and the same statistic give the
opposite: variance roughly twice as high in half-occlusions, which is the safe
direction.

So the failure is not a stimulus-independent property of stereo confidence. It
appears on real imagery and not on synthetic dots, and whether that difference
is about realism or about contrast is unresolved --- a random-dot field is
contrast-uniform by construction and cannot separate them. That open question is
the sharper motivation: confidence fails where the scene is real, and the
synthetic fixture this project used sits on the safe side of the reversal.

\subsection{Concepts}

Four distinctions do most of the analytical work, and none was fully in place
when the project began.

\textbf{Allocation versus acquisition.} \emph{Where to spend limited processing
on data already sampled}, against \emph{where to point the sensor to obtain data
not yet sampled}. Most treatments of active vision conflate them. Separating
them is what makes the results here interpretable, and the separation has a
clean geometric form: a rotation about a fixed optical centre changes allocation
and cannot change acquisition.

\textbf{Unknowable is not uncertain.} A pixel with no valid correspondent is not
a pixel we are unsure about; it is one we cannot measure. The difference between
a binary \emph{can this be measured} and a continuous \emph{how well} turned out
to dominate.

\textbf{Foveal weighting versus foveal sampling.} Weighting says the fovea
contributes more to the posterior. Sampling says the periphery was never
resolved in the first place. On a uniformly sampled sensor the first can only
discard information already paid for. Only the second is what a retina does.

\textbf{Linearising about the fixation is a confidence mechanism.} Expanding the
depth relation about the fixated disparity is exact at the fovea and degrades
with the square of the offset. That degradation term does more than bound
approximation error: it inflates the variance of any estimate far from the
fixated depth, and thereby refuses precisely the confident-and-wrong matches
described above.

A fifth concept was not brought to the project but forced on it, and \S2 begins
there: a stimulus is an instrument, and must be characterised before it is used
to measure.

\subsection{Architecture}

\begin{figure}[htbp]
  \centering
  \includegraphics[width=\textwidth]{architecture.pdf}
  \caption{The six-layer framework, with the three components that sit outside
  it and the three that are named but not built. \emph{Measured} means the
  ledger contains a measurement of the layer's behaviour, and every such box
  names the foreclosure or measurement entry it rests on; \emph{built} means
  implemented and tested with no measurement bearing on the question the layer
  exists to answer; \emph{hypothetical} means recorded as an open decision and
  not built. Drawn from the structure of \texttt{docs/state.yaml} and
  \texttt{docs/inherited-measurements.yaml}, and the script asserts that every
  cited id exists in one of them before it draws. The layer numbering
  \texttt{L1}--\texttt{L6} is Chat-surface vocabulary and is not defined in the
  repository (\texttt{fc-011}).}
  \label{fig:architecture}
\end{figure}

The framework is six layers, and a pass through the loop is the clearest way to
describe them.

A fixation is a gaze direction and a vergence. \textbf{L1} turns it into two eye
rotations under Listing's law, which places the cameras. \textbf{L2} encodes the
resulting pair; \textbf{L3} matches, producing disparity with a variance and a
validity mask. \textbf{L4} scales disparity to metric depth by expanding about
the fixated disparity --- first-order-exact at the fovea, degrading as the
square of the offset. \textbf{L5} decides where attention is allocated;
\textbf{L6} picks the next fixation from the belief's uncertainty. The belief
updates, and the loop repeats.

The layers are not equally weighted in practice. L1 is substantial and was
expensive to get right; L2 and L3 are deliberately minimal; L4 carries more than
its size suggests, through the Jacobian that converts disparity variance into
depth variance; L5 was decided as attentional allocation with the mechanical
branch left untested; L6 is a two-line argmax that, as \S2.2 reports, turned out
not to matter.

Three components sit outside the six-layer stack, and each exists because the
stack alone was insufficient.

\textbf{The sampling model} maps a sample index to a unit ray. It is what makes
the projection swappable: a spherical sensor satisfied the same interface as a
pinhole one without a signature change, because the frame is identical and only
the index-to-direction map differs --- which is the whole of what a projection
is.

\textbf{The belief} is head-anchored and stores directions rather than pixels,
so it survives a change of projection untouched. Its failure mode was
instructive: tied initially to the rectifying rotation, which ignores azimuth by
construction, it silently fused measurements of different world directions into
the same cell, and no test caught it because every control varied elevation.

\textbf{The rig/sensor split.} The original rig type carried baseline and focal
length together, which coupled two thirds of the geometry tests to an intrinsic
they did not need. Separating stereo geometry from the sensor is what made the
rotational content portable.

Listing's tilt is a parameter, not a constant: the default is 0.25 per radian of
vergence, while 0.5 is the value at which vertical disparity in the plane of
regard vanishes exactly. The two are distinct and the distinction is pinned by
test.
